\documentclass[11pt,a4paper]{article}
\usepackage[margin=25mm]{geometry}
\usepackage{amsmath,amssymb,amsthm,hyperref}
\newtheorem{lemma}{Lemma}
\newcommand{\gen}[1]{\langle#1\rangle}
\title{The marked two-swap saturation cases}
\author{Human Compression project: bounded L044 proof component}
\date{}
\begin{document}\maketitle
Use the accepted block-deficit setting: $G\leq A_n$ is a
transitive string C-group of degree $n\geq12$, every maximal omission is
intransitive, and there is a unique eligible $i$-edge with
endpoints $a,b$ in the two $G_i$-orbits $A,B$. Suppose both
orbits have two blocks, of sizes $k,l\geq2$, and
$\alpha=\beta=0$. Put $K=k+l$ and $s=\rho_i$.
The accepted saturation lemma gives even $k,l$ and a
faithful $M\cong C_2^2$, acting on each orbit as the diagonal
two-block image. We preserve the literal distinguished
generators; no product replacement is made.

\begin{lemma}[The marked basis is pure]
After interchanging $A,B$ if necessary,
$M=\{\tau_A,\tau_B\}$ with
$\tau_A=\rho_{i-1}$ exchanging the blocks of $A$ and fixing
$B$ pointwise, and $\tau_B=\rho_{i+1}$ doing the reverse.
Both $x=\rho_{i-2}$ and $y=\rho_{i+2}$ exist.
\end{lemma}
\begin{proof}
Every generator outside $i-1,i+1$ fixes both crossing
endpoints, so its images on the two two-block systems are
trivial. The two independent images in $M$ must therefore
occur at these two neighboring indices. They commute.

Suppose both neighboring generators exchange the blocks of
$A$. Since $\beta=0$, they act identically there: each swaps
the two points in every row. Let $a'$ be the point opposite
$a$. Every other generator except possibly $x,y$ commutes
with both swaps, and so fixes $a'$ as well as $a$.
The generator $x$, if present, commutes with $\rho_{i+1}$,
and $y$, if present, commutes with $\rho_{i-1}$; each also
fixes $a'$. Thus all of $G_i$ preserve $\{a,a'\}$,
contradicting transitivity on $A$ because $k\geq2$.
The same argument applies to $B$. Each block character is
therefore nontrivial on exactly one of the two generators;
independence of their simultaneous images makes those choices
different. This is precisely the asserted pure basis in the
original ordering.

Every generator other than $x,\tau_A$ fixes $a'$: for
$y$ and all distant generators use commutation with
$\tau_A$, while $\tau_B$ is trivial on $A$.
If $x$ were absent, the pair $\{a,a'\}$ would again be
invariant. Hence $x$ exists, and must move $a'$ to another
row. Symmetrically $y$ exists and moves the opposite point
$b'$ to another row of $B$.
\end{proof}

\begin{lemma}[Two-row orbits are excluded]
In this setting neither $k$ nor $l$ can equal two.
\end{lemma}
\begin{proof}
Suppose $k=2$. The generator $x$ fixes $a$ and preserves
each of the two blocks of $A$, so is the identity on the
first block. Since it must move $a'$, it is exactly the
single transposition on the other block. On $B$, the
generator $x$ commutes with $\tau_B$ and preserves the
two blocks, so acts diagonally and has even sign. Thus
$x$ is odd on all points, a contradiction. The proof for
$l=2$ is symmetric.
\end{proof}

Let $C$ be the generators outside $x,\tau_A,s,\tau_B,y$.
They are diagonal on the row pairs of both original orbits,
and fix both row anchors. Their row action is faithful,
so $\gen C\leq S_{k-1}\times S_{l-1}$, with its full
independence and IP retained.

\begin{lemma}[The fully saturated rank is impossible]
One cannot have $\gamma=0$, or equivalently $r=K+1$.
\end{lemma}
\begin{proof}
By the preceding lemma take $k,l\geq4$. With $\gamma=0$,
the family $C$ has $K-4$ generators, the maximal independent
length in $S_{k-1}\times S_{l-1}$. The accepted equality
clause identifies its group as that full product.
Its prefix $P$ at indices at most $i-3$ and suffix $Q$
at indices at least $i+3$ commute and have trivial
intersection by IP. Since the two symmetric factors are
centerless and directly indecomposable, the interval groups
are products of the whole factors.

The first row factor must occur in $P$. Otherwise $x$
centralizes its full natural $S_{k-1}$, whose action on
each $k$-point block has a unique fixed point and trivial
centralizer. This would make $x$ fix the opposite anchor
$a'$, a contradiction. Likewise the second factor must
occur in $Q$. Thus $P=S_{k-1}$ and $Q=S_{l-1}$,
each with a maximal independent string marking.

The generator $x$ fixes $B$ pointwise: its row action there
centralizes the natural $Q$ and fixes its unique row anchor.
If $k\geq6$, the ordinary Coxeter marking of $S_{k-1}$
has far-end omission $S_{k-2}$ fixing the anchor and one
other row. There is no outer degree-six exception because
$k$ is even. Since $x$ centralizes that omission and fixes
$a$, it is the identity on the first block of $A$ and
the single transposition involving $a'$ on the other.
It is odd, impossible. This uses the accepted L047
maximal-string classification, including degree five.
Symmetrically $l\geq6$ is impossible.

The only case left is $k=l=4$. On $A$, the prefix has
generators $(1\ 2)(1'\ 2')$ and $(2\ 3)(2'\ 3')$ on
the three nonanchor rows. The generator $x$ centralizes
the first and fixes $0=a$, while it must move $0'=a'$.
Its evenness and the fact that it fixes $B$ give precisely
\[
 (1\ 2)(0'\ 3')\quad\hbox{or}\quad
 (1'\ 2')(0'\ 3').
\]
As proved explicitly in the accepted four-row endpoint,
these three generators have natural symmetric images on
the three first-block nonanchor points and all four
second-block points, and their centralizer fixes those
seven points individually. The generator $s$ commutes
with them, so fixes $A\setminus\{a\}$ pointwise.
The symmetric argument with $Q,y$ shows that it fixes
$B\setminus\{b\}$ pointwise. Consequently $s=(a\ b)$
is odd, the final contradiction.
\end{proof}

\paragraph{Remaining scope.}
This component proves the pure-basis reduction for
$\beta=0$ without changing the marking and excludes
$\gamma=0$. The case $\beta=0$, $\gamma=1$ leaves a
faithful independent string of length $K-5$ in
$S_{k-1}\times S_{l-1}$, with $k,l\geq4$ even; its
classification is not asserted here. The case $\beta=1$
has a regular four-point $M$-orbit and is also separate.
All other outstanding L044 obligations remain pending.
This is not a whole-source recovery or a new question claim.
\paragraph{Attribution and paid dependencies.}
The source problem and orbit-cut framework are from Cameron, Fernandes,
Leemans and Mixer, \emph{Highest rank of a polytope for $A_n$},
arXiv:1605.09173v1, Section~5.1. The present marked-basis and endpoint
arguments are new proofs of bounded cases in that framework. They use
the complete accepted block-deficit, small-endpoint and two-row notes,
and the full retained L047 extremal/normal-quotient interfaces, including
Whiston's independent-generation bounds. Their exact files and accepted
reviews are bound in the companion dependency map; no sibling proof is
an uncharged premise. No new finite census or formal theorem is used.
\end{document}
